\chapter{Reductio ad Absurdum}
\index{Reductio ad Absurdum}
\section{What is this method?}
\emph{Reductio ad absurdum} means ``reduction to absurdity.'' In plain language:
\begin{enumerate}
  \item Assume the opposite of what you want to prove.
  \item Follow the logic carefully.
  \item Reach an impossibility (a contradiction).
  \item Conclude your opposite assumption was wrong.
  \item Therefore, the original statement must be true.
\end{enumerate}

This method is also called \emph{proof by contradiction}.

\section{Euclid's classic example: infinitely many primes}
\index{Infinitely many primes}
We will use this method to prove:
\begin{quote}
There are infinitely many prime numbers.
\end{quote}

\subsection*{Step 1: Assume the opposite}
Assume there are only finitely many primes. Then we can list all of them:
\[
p_1,\ p_2,\ p_3,\ \ldots,\ p_n.
\]
This list is supposed to contain \emph{every} prime.

\subsection*{Step 2: Build a new number}
Now make the number
\[
N=p_1p_2p_3\cdots p_n+1.
\]
So \(N\) is one more than the product of all primes in the list.

\subsection*{Step 3: Check divisibility}
Take any prime \(p_i\) from the list. Since \(p_i\) divides the product
\(p_1p_2\cdots p_n\), dividing \(N\) by \(p_i\) leaves remainder \(1\), not \(0\).

So none of \(p_1,p_2,\ldots,p_n\) divides \(N\).

\subsection*{Step 4: Contradiction}
Every integer greater than \(1\) is either prime or has a prime factor.
\begin{itemize}
  \item If \(N\) is prime, then it is a prime not in our ``complete'' list.
  \item If \(N\) is composite, it has a prime factor, but that prime factor is
  not in the list either (because none of the listed primes divides \(N\)).
\end{itemize}

Either way, we get a new prime outside the list. This contradicts our assumption
that the list already had all primes.

\subsection*{Step 5: Conclusion}
Our assumption (``only finitely many primes exist'') is false. Therefore:
\[
\text{There are infinitely many prime numbers.}
\]

\section{Why this proof is beautiful}
The proof is short, but very powerful:
\begin{itemize}
  \item It does not try to find all primes.
  \item It shows that any finite list of primes can always be beaten.
  \item So the process never ends.
\end{itemize}

\section{15 Puzzle: Sam Loyd's Offer}
Loyd's offer of a $1,000 prize (equivalent to $35,833 in 2025) to anyone 
who could provide a solution for achieving a particular combination specified by Loyd, 
namely reversing the 14 and 15. 

We now prove, by reductio ad absurdum, that this target position is impossible.

Assume for contradiction that it \emph{is} possible, and that it can be reached
in \(n\) moves.

On the one hand, every move changes the empty cell by one square: either one
row up/down or one column left/right. In Loyd's start and target positions, the
empty cell is in the same place, so after \(n\) moves it must have made an even
number of such one-step changes. Therefore \(n\) is even.

On the other hand, label the empty cell as tile \(16\). Each move is then a
single swap between tile \(16\) and one neighboring tile. So each move performs
exactly one swap. From the start position to Loyd's target position, only tiles
\(14\) and \(15\) are exchanged, which is one swap in total. Therefore the total
number of swaps---hence the number of moves \(n\)---must be odd.

So \(n\) would have to be both even and odd, which is impossible. This
contradiction shows that Loyd's target position cannot be solved. \index{15 Puzzle}
